\chapter{Memória}

\chapterauthor{Felipe da Rocha, Felipe Ferreira e Pedro Augusto}

%-------------------------------------------------------------------------------
\section{Descrição}
%-------------------------------------------------------------------------------
% Descreva em alto nível o que este módulo faz e qual problema ele resolve
% dentro do sistema operacional.
%
O módulo de memória é responsável pelo gerenciamento da memória física e virtual do sistema operacional. Ele fornece mecanismos para controle das páginas disponíveis na RAM, alocação de memória para outros módulos do kernel e configuração da Unidade de Gerenciamento de Memória (MMU).

A implementação foi dividida em três componentes principais:

\begin{itemize}
\item \textbf{PMM (Physical Memory Manager):} responsável pelo gerenciamento da memória física através de um bitmap.
\item \textbf{VMM (Virtual Memory Manager):} responsável pela criação das tabelas de páginas e ativação da MMU.
\item \textbf{API de Memória:} interface utilizada pelos demais módulos para solicitar e liberar memória.
\end{itemize}

Essa divisão permite separar responsabilidades e simplificar futuras expansões do sistema.

%-------------------------------------------------------------------------------
\section{Decisões de Arquitetura e Design}
%-------------------------------------------------------------------------------
% Justificar o porque de ter feito as coisas de uma certa maneira.
%
% Tópicos a cobrir:
% - Quais alternativas foram consideradas?
% - Por que a abordagem escolhida foi a melhor para este projeto?
% - Quais são las limitações ou trade-offs da abordagem escolhida?
%

Para o gerenciamento da memória física foi escolhido o uso de um \textbf{bitmap de páginas}. Nesse modelo, cada bit representa uma página de memória de 4 KB:

\begin{itemize}
\item Bit 0: página livre.
\item Bit 1: página ocupada.
\end{itemize}

Outras alternativas consideradas foram listas encadeadas de blocos livres. Entretanto, o bitmap foi escolhido devido à sua simplicidade de implementação e ao baixo consumo de memória, características importantes para um kernel educacional.

O bitmap é armazenado logo após o final da imagem do kernel na memória. Durante a inicialização, todas as páginas são marcadas como ocupadas e, posteriormente, as regiões livres são liberadas.

Para o gerenciamento virtual foi adotado o esquema de paginação, composto por:

\begin{itemize}
\item Diretório de páginas de primeiro nível (L1);
\item Tabelas de páginas de segundo nível (L2);
\item Mapeamento identidade (virtual = físico) durante a inicialização.
\end{itemize}

O mapeamento identidade foi utilizado para simplificar a ativação da MMU, garantindo que o kernel continue executando normalmente após a mudança para memória virtual.

Como limitação, o sistema ainda não implementa paginação por demanda, swapping ou algoritmos avançados de substituição de páginas.


%-------------------------------------------------------------------------------
\section{Detalhes da Implementação}
%-------------------------------------------------------------------------------
% Descreva como a solução foi implementada, conectando o código
% à teoria.
%
% Tópicos a cobrir:
% - Principais estruturas de dados (structs) e suas funções.
% - Algoritmos chave utilizados.
% - Fluxo de execução das funções mais importantes.
% - Inclua pequenos trechos de código com o pacote 'listings'.
%

\subsection{Gerenciador de Memória Física (PMM)}

A estrutura principal utilizada é:

\begin{lstlisting}[language=C]
typedef struct {
uint32_t *map;
uint32_t size;
} Bitmap;
\end{lstlisting}

Onde:

\begin{itemize}
\item \texttt{map}: vetor contendo os bits que representam as páginas.
\item \texttt{size}: quantidade de inteiros utilizados pelo bitmap.
\end{itemize}

\subsubsection{Inicialização}

A função \texttt{pmm\_init()} configura o bitmap e libera as páginas disponíveis após o kernel.

Fluxo de execução:

\begin{enumerate}
\item Posiciona o bitmap após o final da imagem do kernel.
\item Calcula o número total de páginas da RAM.
\item Marca todas as páginas como ocupadas.
\item Calcula o primeiro endereço livre.
\item Libera todas as páginas até o final da RAM.
\end{enumerate}

\subsubsection{Liberação de Página}

\begin{lstlisting}[language=C]
void pmm_free_page(uintptr_t addr)
{
...
bitmap.map[idx] &= ~mask;
}
\end{lstlisting}

A operação realiza um AND com a máscara invertida, fazendo com que o bit correspondente seja zerado.

\subsubsection{Alocação de Página}

\begin{lstlisting}[language=C]
void* pmm_alloc_block()
{
...
    bitmap.map[i] |= mask;
    return (void*)phys_addr;
}
\end{lstlisting}

O algoritmo percorre o bitmap procurando o primeiro bit livre e o marca como ocupado.

\subsection{Gerenciador de Memória Virtual (VMM)}

A função \texttt{vmm\_init()} realiza a configuração da memória virtual.

As etapas executadas são:

\begin{enumerate}
\item Inicialização do diretório de páginas.
\item Mapeamento da UART.
\item Mapeamento identidade de toda a RAM.
\item Carregamento do diretório no registrador TTBR0.
\item Ativação da MMU.
\end{enumerate}

\subsubsection{Mapeamento de Página}

A função \texttt{vmm\_map\_page()} cria ou reutiliza uma tabela L2 e insere uma entrada apontando para a página física desejada.

\begin{lstlisting}[language=C]
pt->entries[pt_idx] =
(physical_addr & 0xFFFFF000)
| flags
| 0x2;
\end{lstlisting}

\subsubsection{Ativação da MMU}

\begin{lstlisting}[language=C]
void vmm_enable()
{
...
}
\end{lstlisting}

A função \texttt{vmm\_enable()} realiza:

\begin{itemize}
\item Configuração do DACR;
\item Leitura do registrador SCTLR;
\item Habilitação do bit de ativação da MMU;
\item Limpeza do pipeline através da instrução ISB.
\end{itemize}

\subsection{API de Memória}

A API fornece uma interface simplificada para os demais módulos.

Atualmente estão implementadas:

\begin{itemize}
\item \texttt{kmalloc()} -- reserva memória.
\item \texttt{kfree()} -- libera memória.
\end{itemize}

A implementação de \texttt{kmalloc()} ainda encontra-se em desenvolvimento e atualmente calcula o número de páginas necessárias para a alocação.


%-------------------------------------------------------------------------------
\section{Interface Pública (API)}
%-------------------------------------------------------------------------------
% Documente como outros módulos do kernel podem usar este módulo.
% Essencialmente, é a documentação do arquivo .h correspondente.
%

\subsection{API de Uso Geral (Gerenciamento de Heap)}
Esta subseção documenta as funções de alto nível voltadas para a alocação dinâmica de memória no kernel. Esta é a interface padrão exposta e visível para a grande maioria dos outros subsistemas do kernel (como sistemas de arquivos, gerência de rede e drivers comuns de dispositivos).

\subsubsection{kmalloc}

\textbf{Descrição:}
Solicita memória dinâmica ao kernel.

\textbf{Parâmetros:}
\begin{itemize}
\item \texttt{size}: quantidade de bytes desejada.
\end{itemize}

\textbf{Retorno:}
Endereço da memória virtual alocada no Heap.

\subsubsection{kfree}

\textbf{Descrição:}
Libera uma área de memória anteriormente reservada no Heap.

\textbf{Parâmetros:}
\begin{itemize}
\item \texttt{addr}: endereço de memória a ser liberado.
\end{itemize}

\textbf{Retorno:}
\begin{itemize}
\item 1 em caso de sucesso.
\end{itemize}


\subsection{API de Baixo Nível / Core (Subsistemas Críticos)}
Esta subseção documenta as rotinas de controle direto sobre a memória física (PMM) e virtual (VMM). O acesso a estas funções é restrito e especializado, sendo direcionado principalmente para o código de boot do sistema (\texttt{kmain}), o gerenciador de processos (scheduler) para criação de novos contextos/espaços de endereçamento, e drivers de baixo nível que exigem mapeamento direto de registradores de hardware (MMIO).

\subsubsection{pmm\_init}

\textbf{Descrição:}
Inicializa o gerenciador de memória física e configura a estrutura interna do bitmap. Deve ser invocada estritamente durante a rotina de boot.

\textbf{Parâmetros:}
Nenhum.

\textbf{Retorno:}
Não possui retorno.

\subsubsection{pmm\_alloc\_block}

\textbf{Descrição:}
Procura e aloca uma página física livre diretamente da RAM.

\textbf{Parâmetros:}
Nenhum.

\textbf{Retorno:}
\begin{itemize}
\item Endereço da página física alocada.
\item NULL caso não exista memória livre.
\end{itemize}

\subsubsection{pmm\_free\_page}

\textbf{Descrição:}
Libera uma página física previamente alocada, devolvendo-a para o bitmap de páginas livres.

\textbf{Parâmetros:}
\begin{itemize}
\item \texttt{addr}: endereço físico da página.
\end{itemize}

\textbf{Retorno:}
Não possui retorno.

\subsubsection{vmm\_init}

\textbf{Descrição:}
Inicializa a estrutura base do sistema de memória virtual e ativa a MMU do processador. Deve ser invocada apenas uma vez no boot.

\textbf{Parâmetros:}
Nenhum.

\textbf{Retorno:}
Não possui retorno.

\subsubsection{vmm\_map\_page}

\textbf{Descrição:}
Mapeia explicitamente uma página virtual para uma página física correspondente no diretório de tabelas. Útil para mapeamento de hardware (MMIO) e isolamento de processos.

\textbf{Parâmetros:}
\begin{itemize}
\item \texttt{dir}: diretório de páginas (L1).
\item \texttt{virtual\_addr}: endereço virtual desejado.
\item \texttt{physical\_addr}: endereço físico de destino.
\item \texttt{flags}: permissões e flags de acesso à página (ex: leitura, escrita, cache).
\end{itemize}

\textbf{Retorno:}
\begin{itemize}
\item 0 em caso de sucesso.
\item -1 em caso de erro.
\end{itemize}


%-------------------------------------------------------------------------------
\section{Testes e Verificação}
%-------------------------------------------------------------------------------
% Mostre como o funcionamento do módulo foi testado e validado.
%
% Tópicos a cobrir:
% - Descreva a rotina de teste que foi adicionada à função kmain().
% - Inclua a saída esperada (e obtida) na console serial do QEMU.
% - Explique quais casos de borda foram testados.
%

Os testes foram executados dentro do ambiente QEMU utilizando a porta serial para exibição das mensagens de depuração.

A rotina de testes foi adicionada à função \texttt{kmain()}.

Durante o boot são executadas as seguintes etapas:

\begin{enumerate}
\item Inicialização da serial.
\item Inicialização do PMM.
\item Teste da API de memória.
\item Inicialização do VMM.
\item Ativação da MMU.
\end{enumerate}

Saída esperada:

\begin{verbatim}
--- KERNEL BOOT: TESTE DE MEMORIA ---
[INIT] Inicializando PMM...
[TEST] Testando kmalloc...
[TEST] kmalloc alocou memoria com sucesso!
[INIT] Inicializando VMM e ativando MMU...
[SUCCESS] MMU ativa e rodando em modo mapeado!
\end{verbatim}

Os seguintes casos foram validados:

\begin{itemize}
\item Inicialização correta do bitmap.
\item Liberação das páginas após o kernel.
\item Alocação de páginas físicas livres.
\item Criação dinâmica de tabelas de páginas.
\item Mapeamento identidade da RAM.
\item Ativação da MMU sem geração de exceções.
\end{itemize}

Os resultados obtidos foram compatíveis com o comportamento esperado, demonstrando o funcionamento correto da infraestrutura básica de gerenciamento de memória do sistema operacional.